\section{Triton kernel 参数表}

\begin{longtable}{p{0.22\textwidth}p{0.28\textwidth}p{0.40\textwidth}}
\toprule
\textbf{参数/概念} & \textbf{出现位置} & \textbf{含义} \\
\midrule
\texttt{BLOCK\_SIZE} & GeLU、softmax、row sum & 一个 program 处理多少元素或列。太小复用差，太大 register/shared memory 压力高。 \\
\texttt{tl.program\_id} & 所有 Triton kernels & 当前 program/block 的 id，用来决定处理哪一块数据。 \\
\texttt{offsets} & GeLU/softmax & 当前 block 对应的全局元素下标。 \\
\texttt{mask} & 尾部处理 & 防止读写越界，是 block size 不整除数据长度时的必要保护。 \\
\texttt{stride} & softmax/matmul & 描述逻辑二维 tensor 在一维内存中的跳步方式。 \\
\texttt{BLOCK\_M/N/K} & matmul & C tile 的行/列和 K 维 chunk 大小，决定 tile shape 和复用。 \\
\texttt{tl.dot} & matmul & 调用 Triton 的 block-level dot，通常映射到底层矩阵硬件。 \\
\bottomrule
\end{longtable}

\begin{knowledgebox}{为什么这些参数不是随便调}
Triton kernel 的参数直接映射到硬件资源：block size 影响并行粒度，mask 影响边界处理，stride 决定内存访问连续性，tile shape 影响 shared memory/register pressure 和 tensor core 使用。调参不是玄学，而是在硬件约束下找平衡。
\end{knowledgebox}



